\chapter{Data and Panel Construction}
\label{ch:data}

This chapter documents the panel the hypothesis chapters draw on. It comes first because the
construction is shared, and because several choices made here --- the year-alignment rule above all
--- silently corrupt results if got wrong and are invisible in the output if got right. Every
choice was logged when it was taken, with its rationale; the log is Appendix~\ref{app:decisions},
and decisions are cited below by identifier (D1, D2, \ldots).

\section{Unit of analysis and coverage}
\label{sec:data-unit}

The unit of observation is the \textbf{state--financial year}, and every variable is a
\textbf{stock as at 31 March}, not a flow. A financial year is identified by the calendar year of
the 31 March on which it closes: \texttt{fy\_end}~$=2018$ denotes the position as at 31 March 2018,
the close of financial year 2017--18.

\begin{table}[H]
\centering
\caption{Panel coverage}
\label{tab:panels}
\small
\begin{tabular}{lccccl}
\toprule
Panel & Window & Years & States & $N$ & Infrastructure measured \\
\midrule
Track 1 & \panelA{} & 6 & 33 & 198 & Branches and ATMs \\
Track 2 & \panelB{} & 12 & 32 & 384 & Branches only \\
Track 3 & FY2019--FY2023 & 5 & 33 & 165 & Branches, ATMs and UPI \\
\bottomrule
\end{tabular}
\end{table}

\noindent
All three panels are \textbf{balanced and complete}: every state appears in every year, with no
missing cells and no non-positive values in any series. The state counts differ for one documented
reason given in Section~\ref{sec:data-harmonisation}. ATMs cannot be extended before 2018, because
RBI's state-wise release begins there, which is why Track~2 carries branches alone; Track~3 starts
at FY2019 because that is the first complete financial year of digital payment data. Track~3's
banking columns are asserted \textbf{identical} to Track~1's, so the earlier results remain
reproducible from an untouched file and the tracks stay directly comparable.

\section{Sources}
\label{sec:data-sources}

Nine series are used. Eight are official publications of the RBI or the Government of India, with
one academic redistribution filling a single year; the ninth is the digital payment series
introduced for Track~3.

\begin{table}[H]
\centering
\caption{Source of each series}
\label{tab:sources}
\scriptsize
\begin{tabularx}{\textwidth}{@{}p{2.7cm}p{3.3cm}Xp{1.7cm}@{}}
\toprule
\textbf{Series} & \textbf{Publication} & \textbf{Table / report} & \textbf{Role} \\
\midrule
Bank branches & RBI \emph{Handbook of Statistics on Indian States}, 2024--25 &
Table 152, bank offices of scheduled commercial banks, end-March & Independent \\
\addlinespace
ATMs & RBI \emph{State-wise and Region-wise Deployment of ATMs} &
Standalone quarterly release; end-March files 2018--2024 & Independent \\
\addlinespace
Deposit accounts & RBI Basic Statistical Returns (BSR-2), via DBIE/CIMS &
Report 853, Annual BSR-2 Table 2.3; 31 March 2018 from India Data Portal & \Access{} \\
\addlinespace
Credit accounts & RBI Basic Statistical Returns (BSR-1), via India Data Portal &
Derived state--year credit accounts 2010--2023, including RRBs & \Access{} \\
\addlinespace
Deposits outstanding & RBI \emph{Handbook} & Table 155, deposits of SCBs (\rs crore) & \Usage{} \\
\addlinespace
Credit outstanding & RBI \emph{Handbook} & Table 156, credit of SCBs (\rs crore) & \Usage{} \\
\addlinespace
Per-capita NSDP & RBI \emph{Handbook} & Table 19, per-capita NSDP, current prices & Control \\
\addlinespace
Adult population & Technical Group on Population Projections (2020) &
Annual totals with quinquennial age structure; 1 March series & Denominator \\
\addlinespace
Digital payments & PhonePe Pulse open data release &
Registered users, transaction count and value, by state and quarter & Independent \\
\bottomrule
\end{tabularx}
\end{table}

\noindent
Retrieval was non-trivial in each case. The Handbook tables \citep{rbi_handbook} are HTML only,
each split across two blocks covering earlier and later years, and were parsed and stitched
programmatically. The ATM release \citep{rbi_atm} is navigable only through an ASP.NET postback
widget rather than plain URLs, so per-release permalinks were harvested by scripting those
postbacks. The BSR returns \citep{rbi_bsr} are no longer published as static files at all --- the
legacy links redirect to a single-page application --- and were retrieved programmatically from the
DBIE/CIMS API; the 31 March 2018 deposit-account figures, which the current report does not reach
back to, come from an academic redistribution of the same returns \citep{idp}. Adult population
derives from the official projections report \citep{popproj2020}.

\section{The year-alignment rule}
\label{sec:data-years}

Three year conventions are in play, and they disagree on labels even when they agree on dates. This
is the easiest place in the project to introduce a silent one-year offset --- an error that
produces no missing value, no warning, and no implausible number, and simply biases every estimate.
The Handbook banking tables and the ATM release label a stock as at 31 March 2018 as
\texttt{2018}; the Handbook NSDP tables call the same financial year \texttt{2017-18}; the
population projections date their stock \texttt{1 March 2018}.

\begin{keybox}{Alignment rule (D5, D9)}
Banking year $Y$ $\;\equiv\;$ NSDP financial year $(Y-1)$--$Y$ $\;\equiv\;$ population at
1 March $Y$.

\medskip
All three are recorded in one canonical column, \texttt{fy\_end}, added to all eight input files.
Merging and filtering use \texttt{fy\_end} only, and it supplies the year fixed effect.
\end{keybox}

\noindent
Of the eight files only the income file required conversion; native labels were retained alongside
\texttt{fy\_end} so the original coding stays auditable. The 1 March population series is used
rather than 1 July or 1 October: every numerator is a stock at 31 March, so the denominator must be
a stock at approximately the same instant --- one month away rather than six. The rule was verified
against published figures rather than assumed, as Section~\ref{sec:data-validation} records.

\section{State harmonisation}
\label{sec:data-harmonisation}

Fixed effects require each unit to denote the same thing in every year. India's boundaries changed
inside both windows and, more awkwardly, different sources record the same change in different
years. Three adjustments follow.

\textbf{Units dropped (D1).} Lakshadweep and the combined Dadra \& Nagar Haveli and Daman \& Diu
are dropped: no per-capita NSDP is published for either in any year, so neither can carry the
income control. Both are very small; the panel goes from 35 units to 33.

\textbf{Units merged (D2).} Jammu \& Kashmir with Ladakh, and Dadra \& Nagar Haveli with Daman and
Diu, \emph{in every year}. Merging only from the year of the formal change fails because the
sources disagree about when that was: in the Handbook, Daman \& Diu goes blank from 2020 exactly as
Dadra \& Nagar Haveli jumps from 59 to 110 branches, whereas the ATM files keep three separate
columns until 2021. Summing all components in all years is the only rule producing a continuous
series in both sources simultaneously.

\textbf{Andhra Pradesh and Telangana, twelve-year panel only (D13).} Telangana was carved out of
Andhra Pradesh in June 2014, so in the twelve-year panel it shows \emph{zero} branches, deposits and
credit for 2012--2014, its activity being recorded inside Andhra Pradesh. Branches confirm the
mechanism exactly: Andhra Pradesh falls from 9,900 in 2014 to 6,290 in 2015 as Telangana appears
with 4,721, and the combined series runs continuously from 9,900 to 11,011. The two are therefore
summed in every year, which is why Track~2 has 32 units. The failure mode is instructive: had the
merge been missed, three state-years would have entered recording zero infrastructure alongside
zero deposits --- outliers positioned precisely to manufacture a spurious \emph{positive}
association between infrastructure and inclusion, the sign the hypothesis predicts. It was caught
by a completeness check flagging zeros, not by inspecting results.

\begin{keybox}{General rule}
Merges are performed on \textbf{raw levels} --- branches, ATMs, deposits, credit, account counts,
population --- \textbf{before} any ratio is computed. Ratios are never averaged across merged units.
\end{keybox}

\noindent
After upper-casing and standardising ``AND'' to ``\&'', 34 of 37 units match exactly across
sources; the exceptions are Delhi, Jammu \& Kashmir, and Daman \& Diu, each named differently in
the population or ATM files. Matching uses \textbf{word boundaries} rather than substrings: a naive
substring test folds \texttt{ANDAMAN} into the Daman group, an error that occurred once during
preparation and is now guarded by an explicit assertion.

\section{Validation}
\label{sec:data-validation}

The panel was checked against independently published aggregates rather than accepted on trust.
All-India ATM totals recomputed from the state columns match RBI's published national figures
exactly in every year 2018--2024 (222,066 in 2018 through 257,654 in 2024; the 2024 decline is
genuine decommissioning, not a parsing error). State-by-year deposit amounts differ from Handbook
Table 155 by a median of 0.0001 per cent, and credit amounts from Table 156 by 0.001 per cent.
Population totals reproduce from the sum of age bands, and the published 18+ figure reproduces from
a residual identity, for all 23 states that print it. On year alignment, Maharashtra at
\texttt{fy\_end}~$=2018$ returns 12,545 branches, 25,651 ATMs and per-capita income of
\rs 172,663 --- all matching RBI's published values exactly, confirming no offset was introduced.
Completeness is asserted in code: 198 and 384 rows, balanced, no nulls, no zeros.

\section{The digital payment series}
\label{sec:data-digital}

Track~3 requires a measure of digital financial infrastructure at state level. India has no
official state-wise UPI release covering the full window, so the series used is PhonePe Pulse
\citep{phonepepulse}, the open data release of one of the two largest UPI applications: registered
users, transaction counts and transaction value, by state and calendar quarter.

Conversion to the panel's year convention follows the same rule as everything else. Transactions
and value are \textbf{flows}, summed over the four quarters making up each financial year;
registered users is a \textbf{stock}, taken at the quarter ending 31 March and verified monotone
across all quarter transitions. Denominators follow the same access/usage distinction as the
banking series: registered users \textbf{per adult}, transactions and value \textbf{per capita}.
Coverage is complete --- 36 raw states across 34 quarters with no gaps and no missing amounts ---
and every one of the 33 panel states is present in every year.

Four independent validations were run rather than taking the release on trust. The transaction-type
split (peer-to-peer, retail, utility) reproduces the aggregate total to \textbf{0.000000 per cent}
across all 1,224 state-quarters, meaning two independently maintained trees within the release
agree exactly. Registered users reach \textbf{45.5 crore} at 31 March 2023, against PhonePe's own
publicly stated figure of over 45 crore. Transactions total \textbf{44.3 billion} in FY2023 against
roughly 84 billion for all of UPI that year \citep{npci2025}, consistent with the firm's
approximately 50 per cent market share. And the implied mean ticket size of \rs 1,593 confirms
amounts are denominated in rupees rather than paise.

Three caveats follow the series into the analysis and are revisited in
Section~\ref{sec:h1-limitations}. \textbf{PhonePe is one firm, not the market}, so this measures
PhonePe penetration as a proxy for digital payment penetration, and because market share varies by
state, cross-state levels are noisier than within-state growth --- which is fortunate, since
within-state growth is what the estimator uses. \textbf{Registered users are accounts, not people},
so one person with two numbers counts twice. And transactions are attributed to the \textbf{user's
registered state}, a third attribution rule alongside the branch-location and place-of-sanction
rules already in the panel.

\section{Known limitations of the data}
\label{sec:data-defects}

Three defects are carried openly rather than silently repaired, and each is revisited where it
bears on results.

\paragraph{A break in RBI's own deposit-account series (D7).} The savings-account column of BSR-2
is discontinuous for three states: Delhi runs 35.5, 38.7, \textbf{106.1}, 130.9, 159.4 million
accounts, breaking from 2023 and persisting; Haryana runs 47.7, 52.6, \textbf{102.1}, breaking from
2024; Chandigarh runs 2.98, \textbf{5.03}, 3.25, spiking in 2023 and reverting. This is an error in
RBI's published source, not in extraction --- the component columns reproduce the published total
exactly wherever both are printed --- and the likely cause is a bank reclassifying savings accounts
to a head-office location. The symptom is visible downstream: deposit accounts per adult reaches
7.99 for Delhi in 2023 and 7.69 for Chandigarh in 2023, against roughly 1 to 3.5 elsewhere. Two of
the four affected cells fall inside \panelA{}. The decision, taken \emph{before} estimation, was to
retain them in the main specification and re-run without them as a robustness check;
Section~\ref{sec:h1-robustness} reports what happened, and it is the most consequential result in
the report. Goa sits at 5.45--5.73 across \emph{all} six years and Chandigarh at 5.1--5.4 outside
2023 --- high but flat, so genuine levels rather than defects.

\paragraph{Adult population is interpolated, not published (D6).} The projections publish annual
totals for 2011--2036 but age structure only at five-year anchors, so there is no published annual
adult count. The 18+ series is \emph{derived}: the adult share is linearly interpolated between
anchors and applied to the published annual total. Totals are official; adult counts are ours. Of
the 266 state-years in the underlying file, 243 rest on an interpolated share. The consequence goes
beyond ordinary measurement error, as Section~\ref{sec:h1-gate} develops: being near-linear, the
series is unusually smooth, and as the denominator of every per-adult ratio it contributes almost no
independent variation --- it deflates a trend without adding movement to it. A further caveat
attaches to a few units: the source excludes Goa from its age-structure exercise and reports the
seven smaller north-eastern states only in combination, so proxy shares were assigned and flagged.

\paragraph{Universe and basis of the credit-account series (D8).} The two halves of \Access{}
initially rested on different bank universes: deposit accounts from BSR-2 \emph{include} regional
rural banks, while the quarterly BSR-1 credit series \emph{excludes} them and records loans by place
of sanction rather than utilisation. Excluding RRBs removes a median of roughly 18 per cent of
credit accounts nationally, but unevenly --- 65 per cent in Mizoram, 39 in Tripura, 33 in Uttar
Pradesh --- and the share drifts within states over time, so state fixed effects do not absorb it.
Critically, it bites hardest in exactly the rural, lower-income states the hypothesis is about. The
series was therefore switched to an RRB-inclusive, place-of-utilisation file covering 2010--2023.
The two could not be reconciled arithmetically --- differencing them produced impossible negative
values for Maharashtra, because they differ on \emph{where} a loan is counted as well as on
\emph{which} banks --- so this is a replacement, not an adjustment. The cost is FY2024, which fixes
\panelA{} at six years; the incidental benefit is that the last missing cell, Gujarat's FY2023--24
income, falls outside the window, leaving the panel perfectly balanced. One asymmetry remains:
deposit accounts are recorded by branch location and credit accounts by place of utilisation, and
neither corresponds to depositor or borrower residence.

\section{What the hypothesis chapters draw from this panel}
\label{sec:data-usage}

Chapter~\ref{ch:h1} uses all nine series: the eight banking and demographic series across all
three tracks, and the digital payment series for Track~3.

\stub{If Hypotheses 2 and 3 use series beyond the eight above, add them to
Table~\ref{tab:sources} and document any extra harmonisation here, so this chapter remains the
single place the panel is described.}
